\chapter{Guillain-Barré Syndrome}

\section*{Definition and Overview}
Guillain-Barré syndrome (GBS) is a rare but serious autoimmune disorder in which the body's immune system attacks the peripheral nerves, the nerves that connect the brain and spinal cord to the rest of the body. This causes progressive weakness, usually beginning in the legs and spreading upward, sometimes affecting the arms, face, and breathing muscles. GBS often follows an infection such as a cold, gastroenteritis, or COVID-19 by days to weeks, and it is thought that the immune response to the infection mistakenly damages the nerves. Most people require hospital care, and treatment with immunoglobulin or plasma exchange speeds recovery. While GBS can be life-threatening, most people recover, though some have lasting weakness.

\section*{Epidemiology}
GBS is rare, affecting around 1--2 people per 100,000 each year, which translates to roughly 300--400 cases in Australia annually. It can occur at any age, with a slight peak in older adults and in young adults, and is slightly more common in men. About two-thirds of cases follow an infection, most commonly campylobacter gastroenteritis, respiratory infections, or the flu. GBS is also a recognised complication of some viral infections, including Zika virus and COVID-19. The incidence increases with age, and it is slightly more common in winter and spring.

\section*{Aetiology and Pathophysiology}
GBS is an autoimmune attack on the peripheral nervous system, often triggered by infection.

\begin{itemize}
    \item \textbf{Trigger:} a preceding infection stimulates an immune response; in susceptible people, antibodies cross-react with components of peripheral nerves (molecular mimicry)
    \item \textbf{Demyelination:} in the most common form (acute inflammatory demyelinating polyneuropathy), the immune system damages the myelin sheath that insulates nerves
    \item \textbf{Axonal damage:} in some forms, the nerve fibres themselves are attacked, causing more severe and slower-recovering weakness
    \item \textbf{Impaired signalling:} damaged nerves cannot transmit signals properly, causing weakness, numbness, and loss of reflexes
    \item \textbf{Respiratory weakness:} the nerves to the breathing muscles can be affected, requiring ventilation
    \item \textbf{Autonomic involvement:} the nerves controlling heart rate and blood pressure can be affected, causing dangerous fluctuations
\end{itemize}

\section*{Clinical Features}
GBS usually follows a characteristic pattern of progressive weakness.

\begin{itemize}
    \item Weakness that begins in the legs and spreads upward to the arms, trunk, and face over days to weeks
    \item Tingling, numbness, or pain in the hands and feet
    \item Loss of reflexes (hyporeflexia)
    \item Difficulty walking, climbing stairs, and rising from a chair
    \item Facial weakness, difficulty swallowing, and double vision
    \item Breathlessness and difficulty breathing in severe cases
    \item Autonomic symptoms: irregular heart rate, blood pressure fluctuations, and bladder problems
    \item Severe pain in the back, limbs, or muscles in many cases
    \item Symptoms typically peak within two to four weeks
\end{itemize}

\section*{Differential Diagnosis}
Other conditions cause progressive weakness and must be excluded.

\begin{itemize}
    \item Chronic inflammatory demyelinating polyneuropathy (CIDP), a related but slower condition
    \item Transverse myelitis and other spinal cord disorders
    \item Botulism and other toxin-induced paralysis
    \item Myasthenia gravis, which causes fluctuating weakness
    \item Tick paralysis
    \item Poliomyelitis and other viral infections affecting nerves
    \item Electrolyte disturbances, including low potassium
    \item Lead poisoning and other neuropathies
    \item Stroke and other central nervous system disease
\end{itemize}

\section*{When to Seek Medical Care}
Seek urgent medical assessment for any progressive weakness, especially weakness spreading upward, tingling that is spreading, or difficulty breathing. GBS is a medical emergency that requires hospital care.

\textbf{Call 000 (or local emergency services) for:}
\begin{itemize}
    \item Difficulty breathing or increasing breathlessness
    \item Rapidly progressive weakness, especially of the arms, legs, or face
    \item Difficulty swallowing or an inability to clear secretions
    \item Collapse, fainting, or severe dizziness from blood pressure problems
\end{itemize}

\section*{Diagnosis and Investigations}
Diagnosis is made from the clinical picture and confirmed with tests.

\begin{itemize}
    \item \textbf{Clinical assessment:} progressive symmetrical weakness with reduced reflexes after an infection
    \item \textbf{Lumbar puncture:} elevated protein in the cerebrospinal fluid without a high cell count supports the diagnosis
    \item \textbf{Nerve conduction studies and EMG:} show slowed or blocked nerve signals and confirm peripheral nerve damage
    \item \textbf{Respiratory monitoring:} measurement of breathing strength (vital capacity) is essential
    \item \textbf{Blood tests:} exclude other causes, including metabolic and toxic causes
    \item \textbf{Antibody tests:} for specific GBS variants in some cases
    \item \textbf{Hospital admission:} all suspected cases are admitted for monitoring and treatment
\end{itemize}

\section*{Management}
GBS is treated in hospital, often in intensive care, with supportive and specific therapies.

\begin{itemize}
    \item \textbf{Intravenous immunoglobulin (IVIG):} the main treatment, given over several days; it blocks the immune attack
    \item \textbf{Plasma exchange (plasmapheresis):} an alternative that filters harmful antibodies from the blood
    \item \textbf{Respiratory support:} ventilation if breathing weakness develops
    \item \textbf{Monitoring:} close observation of breathing, heart rhythm, and blood pressure
    \item \textbf{Pain management:} medications for the nerve pain common in GBS
    \item \textbf{Prevention of complications:} physiotherapy to prevent contractures, blood clot prevention, and care for pressure areas
    \item \textbf{Swallowing support:} tube feeding if swallowing is affected
    \item \textbf{Rehabilitation:} a long program of physiotherapy, occupational therapy, and speech therapy
    \item \textbf{Psychological support:} for the emotional impact of a sudden, severe illness
\end{itemize}

\section*{Complications}
\begin{itemize}
    \item Respiratory failure requiring ventilation
    \item Autonomic dysfunction, including dangerous blood pressure and heart rhythm problems
    \item Deep vein thrombosis and pulmonary embolism from immobility
    \item Infections, including pneumonia and urinary tract infections
    \item Pressure injuries from prolonged bed rest
    \item Severe and chronic pain
    \item Residual weakness, fatigue, and sensory disturbance that can persist for years
    \item In severe cases, permanent disability or death, though this is uncommon with modern care
\end{itemize}

\section*{Prognosis and Outcomes}
Most people with GBS recover, but recovery is slow. The illness peaks within two to four weeks, and improvement typically begins after that, continuing for months to up to two years. Around 80\% of people walk again independently, though about 20\% have lasting weakness or fatigue. About 3--5\% of people die from complications, mainly in older people with severe disease. Treatment with IVIG or plasma exchange speeds recovery and reduces long-term disability. Most people eventually return to work and daily activities, though fatigue may persist.

\section*{Prevention and Self-Care}
\begin{itemize}
    \item There is no proven way to prevent GBS, and it cannot be predicted after infection
    \item Seek urgent medical care for progressive weakness or tingling after an illness
    \item During hospital care, participate actively in physiotherapy and rehabilitation
    \item Follow the rehabilitation program after discharge, as recovery takes months
    \item Be prepared for fatigue and pace activities during recovery
    \item Seek pain management for nerve pain
    \item Access support groups for people with GBS and their families
    \item Attend follow-up appointments with neurologists and rehabilitation specialists
    \item People who have had GBS may be advised about specific vaccine considerations, discussed with their doctor
\end{itemize}

\section*{Sources and Further Reading}
The following reputable sources provide further information for healthcare students and patients seeking reliable, current advice about Guillain-Barré syndrome.

\begin{itemize}
    \item GBS/CIDP Foundation of Australia. \textit{Information and support for GBS.} https://www.gbs-cidp.org.au/
    \item Healthdirect Australia. \textit{Guillain-Barré syndrome.} https://www.healthdirect.gov.au/
    \item Australian Government Department of Health and Aged Care. \textit{Guillain-Barré syndrome information.} https://www.health.gov.au/
    \item NHS. \textit{Guillain-Barré syndrome.} https://www.nhs.uk/
    \item Mayo Clinic. \textit{Guillain-Barré syndrome: symptoms and causes.} https://www.mayoclinic.org/
    \item National Institute of Neurological Disorders and Stroke. \textit{Guillain-Barré syndrome.} https://www.ninds.nih.gov/
\end{itemize}
